\chapter{Eliminasi \texorpdfstring{$\beta$}{Beta}}\label{ch:b1:elimination}

Reaksikan 2-bromo-2-metilpropana dengan natrium etoksida dalam etanol
pada suhu kamar, maka produk utamanya sebuah eter; panaskan campuran yang
sama, maka produk utamanya sebuah alkena, 2-metilpropena, tanpa gugus
etoksi di dalamnya. Ion etoksida tidak berperan sebagai \omterm{def:b1:electronic-effects:nucleophile}{nukleofil} yang
menyerang karbon, melainkan sebagai basa yang mengambil proton dari
karbon di sebelahnya sementara bromidanya lepas. Persaingan antara
substitusi dan eliminasi ini menjalar ke seluruh sintesis organik. Bab
ini memerikan kedua mekanisme eliminasi, syarat geometrisnya yang
istimewa, aturan yang meramalkan alkena mana yang terbentuk, dan cara
mengarahkan sebuah reaksi ke salah satu jalan.

\begin{recall}
SN1 dan SN2 (\cref{ch:b1:nucleophilic-substitution}); \omterm{def:b1:stereochemistry-in-depth:representations}{proyeksi Newman},
\omterm{def:b1:stereochemistry-in-depth:isomers}{konformasi} sikloheksana, dan deskriptor $Z$/$E$
(\cref{ch:b1:stereochemistry-in-depth}); \omterm{def:b1:electronic-effects:intermediates}{karbokation} dan basa
(\cref{ch:b1:electronic-effects}).
\end{recall}

\section{Mekanisme E2}

\begin{definition}[Eliminasi \texorpdfstring{$\beta$}{beta}, E2 dan E1]\label{def:b1:elimination:elimination}
\emph{Eliminasi $\beta$}\index{eliminasi $\beta$} melepaskan \omterm{def:b1:electronic-effects:nucleophile}{gugus pergi} Y
dari sebuah karbon (karbon $\alpha$) dan sebuah proton dari karbon
tetangganya (karbon $\beta$), sehingga terbentuk ikatan rangkap dua di
antara keduanya. Dalam \emph{mekanisme E2}\index{mekanisme E2} sebuah basa
mengambil proton sementara Y lepas, dalam satu \omterm{def:b1:elementary-steps:elementary-step}{tahap elementer}; dalam
\emph{mekanisme E1}\index{mekanisme E1} Y lepas lebih dahulu,
menghasilkan \omterm{def:b1:electronic-effects:intermediates}{karbokation}, yang lalu kehilangan sebuah proton $\beta$.
\end{definition}

\begin{proposition}[Hukum laju E2]\label{prop:b1:elimination:e2-law}
Untuk E2, $v = k[\mathrm{RY}][\mathrm{B}]$, dengan B basanya.
\end{proposition}

\begin{proof}
Satu \omterm{def:b1:elementary-steps:elementary-step}{tahap elementer} bimolekuler, yang melibatkan \omterm{def:b1:nucleophilic-substitution:substitution}{substrat} dan basa:
hukum aksi massa untuk \omterm{def:b1:elementary-steps:elementary-step}{tahap elementer} (\cref{ch:b1:elementary-steps}).
\end{proof}

\begin{definition}[Anti-periplanar dan sin-periplanar]\label{def:b1:elimination:periplanar}
Dua ikatan pada atom-atom yang bertetangga disebut
\emph{anti-periplanar}\index{anti-periplanar} jika terletak pada satu
bidang di sisi yang berlawanan (\omterm{def:b1:stereochemistry-in-depth:conformations}{sudut torsi} $180^\circ$), dan
\emph{sin-periplanar}\index{sin-periplanar} jika terletak pada satu
bidang di sisi yang sama ($0^\circ$).
\end{definition}

\begin{omfigure}
\begin{tikzpicture}[font=\small, >={Stealth}]
  \pic (n) at (0,0) {omnewman={90}{270}};
  \node at (n-f1) {H}; \node at (n-f2) {\ce{CH3}}; \node at (n-f3) {H};
  \node at (n-b1) {Br}; \node at (n-b2) {\ce{CH3}}; \node at (n-b3) {H};
  \node (b) at (-2.6,2.0) {\ce{EtO-}};
  \draw[->, thick, omThm] (b) to[bend left=20] ($(n-f1)+(-0.15,0.1)$);
  \draw[->, thick, omThm] (0.12,0.8) to[bend left=40] (0.12,0.25);
  \draw[->, thick, omThm] (0.12,-0.6) to[bend right=35] ($(n-b1)+(0.25,0.15)$);
  \node[below] at (0,-2.0) {H dan Br anti-periplanar};
  \draw[->, thick] (2.4,0) -- (3.6,0) node[midway, above] {E2};
  \node at (5.6,0) {\chemfig{C(-[:150]H_3C)(-[:210]H)=C(-[:30]H)(-[:-30]CH_3)}};
  \node at (5.6,-0.9) {\cip{E}-but-2-ena};
\end{tikzpicture}

{\small Reaksi E2 2-bromobutana, dipandang sepanjang ikatan C3--C2 (C3 di
depan, C2 yang membawa Br di belakang) pada salah satu dari kedua
\omterm{def:b1:stereochemistry-in-depth:isomers}{konformasi} \omterm{def:b1:elimination:periplanar}{anti-periplanarnya}; \omterm{def:b1:stereochemistry-in-depth:isomers}{konformasi} yang lain, dengan hidrogen
lain pada C3 anti terhadap Br, menghasilkan \cip{Z}-but-2-ena. Tiga
pasang elektron berpindah bersamaan: basa mengambil proton, pasangan
C--H menjadi ikatan $\pi$, pasangan C--Br pergi bersama brom. Gugus-gugus
yang saling berhadapan pada \omterm{def:b1:stereochemistry-in-depth:representations}{proyeksi Newman} berakhir di sisi yang sama
ikatan rangkap dua.}
\end{omfigure}

\begin{proposition}[E2 bersifat anti dan stereospesifik]\label{prop:b1:elimination:anti}
E2 berlangsung dari \omterm{def:b1:stereochemistry-in-depth:isomers}{konformasi} yang ikatan C--H dan C--Y-nya
\omterm{def:b1:elimination:periplanar}{anti-periplanar}. Reaksi itu stereospesifik: \omterm{def:b1:nucleophilic-substitution:substitution}{substrat-substrat} yang
merupakan \omterm{def:b1:stereochemistry-in-depth:enantiomers}{diastereomer} menghasilkan \omterm{def:b1:stereochemistry-in-depth:isomers}{stereoisomer} alkena yang berbeda,
yang ditentukan oleh susunan anti itu.
\end{proposition}

\begin{proof}
Pada \omterm{def:b1:elementary-steps:energy-profile}{keadaan transisi}, \omterm{def:b1:lewis-resonance-vsepr:lewis}{pasangan elektron ikatan} C--H harus memasuki
daerah antiikatan ikatan C--Y sementara kedua karbon menjadi trigonal;
tumpang tindihnya paling baik jika kedua ikatan sejajar, dalam satu
bidang. Susunan anti adalah susunan goyang (berenergi rendah) dan menjaga
basa jauh dari \omterm{def:b1:electronic-effects:nucleophile}{gugus pergi}; susunan sin adalah susunan eklips. Dengan H
dan Y anti, kedua gugus lain pada setiap karbon menjadi tetap: gugus-gugus
yang berada di sisi yang sama pada \omterm{def:b1:stereochemistry-in-depth:representations}{proyeksi Newman} berakhir \textit{cis}
pada ikatan rangkap dua. Sebuah \omterm{def:b1:stereochemistry-in-depth:enantiomers}{diastereomer} mempunyai dua gugus yang
dipertukarkan pada satu karbon, dan menghasilkan alkena yang lain.
\end{proof}

\begin{proposition}[E2 pada sikloheksana]\label{prop:b1:elimination:cyclohexane-e2}
Pada cincin sikloheksana, E2 menuntut \omterm{def:b1:electronic-effects:nucleophile}{gugus pergi} dan hidrogen $\beta$
sama-sama \omterm{def:b1:stereochemistry-in-depth:conformations}{aksial} (\textit{trans}-diaksial); \omterm{def:b1:electronic-effects:nucleophile}{gugus pergi} yang tertahan
\omterm{def:b1:stereochemistry-in-depth:conformations}{ekuatorial} tidak dapat bereaksi sampai cincinnya terbalik.
\end{proposition}

\begin{proof}
Pada karbon-karbon \omterm{def:b1:stereochemistry-in-depth:conformations}{kursi} yang bersebelahan, dua ikatan \omterm{def:b1:stereochemistry-in-depth:conformations}{aksial} bersifat
\omterm{def:b1:elimination:periplanar}{anti-periplanar} (satu ke atas, satu ke bawah, sejajar dengan sumbu);
ikatan \omterm{def:b1:stereochemistry-in-depth:conformations}{ekuatorial} hanya anti terhadap ikatan C--C cincin, tidak pernah
terhadap ikatan C--H.
\end{proof}

\begin{omfigure}
\begin{tikzpicture}[font=\small]
  \pic (a) at (0,0) {omchair};
  \draw[thick, omThm] (a-c1) -- (a-a1) node[above] {Cl};
  \draw[thick] (a-c1) -- (a-e1) node[right] {H};
  \draw[thick, omDef] (a-c2) -- (a-a2) node[below] {H};
  \draw[thick, omDef] (a-c6) -- (a-a6) node[below] {H};
  \node[align=left, anchor=west] at (3.0,0.2) {Cl aksial (ke atas) pada C1;\\H aksial (ke bawah) pada kedua tetangga:\\dua H anti-periplanar};
\end{tikzpicture}

{\small Susunan \textit{trans}-diaksial pada \omterm{def:b1:stereochemistry-in-depth:conformations}{kursi}: klor yang \omterm{def:b1:stereochemistry-in-depth:conformations}{aksial}
\omterm{def:b1:elimination:periplanar}{anti-periplanar} terhadap hidrogen-hidrogen \omterm{def:b1:stereochemistry-in-depth:conformations}{aksial} pada kedua karbon
tetangganya (biru).}
\end{omfigure}

\section{Mekanisme E1}

\begin{proposition}[Hukum laju E1]\label{prop:b1:elimination:e1-law}
Untuk E1, $v = k[\mathrm{RY}]$, dengan laju sama dengan laju ionisasinya.
\end{proposition}

\begin{proof}
Seperti pada SN1 (\cref{prop:b1:nucleophilic-substitution:sn1-law}):
\omterm{def:b1:electronic-effects:intermediates}{karbokation} terbentuk dalam tahap yang lambat dan terpakai dengan cepat,
di sini melalui pelepasan proton kepada pelarut atau \omterm{def:b1:predominant-reaction:strong-acid}{basa lemah};
\omterm{def:b1:elementary-steps:qssa}{pendekatan keadaan tunak} menyisakan $v = k_1[\mathrm{RY}]$.
\end{proof}

E1 berbagi \omterm{def:b1:electronic-effects:intermediates}{karbokationnya} dengan SN1: \omterm{def:b1:nucleophilic-substitution:substitution}{substrat} tersier dalam \omterm{def:b1:intermolecular-forces-solvents:solvent-classes}{pelarut
protik} menghasilkan keduanya, dengan alkena yang lebih disukai jika
dipanaskan. Karena \omterm{def:b1:electronic-effects:intermediates}{karbokation} itu datar dan protonnya baru lepas
sesudahnya, E1 tidak mempunyai syarat geometris dan tidak stereospesifik;
reaksi itu terutama menghasilkan alkena yang lebih stabil, biasanya
isomer \cip{E}.

\begin{omfigure}
\setchemfig{atom sep=2.1em}
\schemestart
\chemfig{C(-[2]CH_3)(-[4]H_3C)(-[6]CH_3)-Br}
\arrow{->[lambat]}
\chemfig{C^{+}(-[2]CH_3)(-[:210]H_3C)(-[:-30]CH_3)}
\arrow{->[cepat][\ce{-H+}]}
\chemfig{C(-[2]CH_3)(-[:210]H_3C)=[:-30]CH_2}
\schemestop
\par\vspace{4pt}

{\small E1: ionisasi 2-bromo-2-metilpropana menjadi \omterm{def:b1:electronic-effects:intermediates}{karbokation}, lalu
pelepasan sebuah proton dari gugus metil kepada pelarut, yang
menghasilkan 2-metilpropena.}
\end{omfigure}

\section{Regioselektivitas: aturan Zaitsev}

\begin{definition}[Reaksi regioselektif dan stereoselektif]\label{def:b1:elimination:selectivity}
Reaksi yang dapat menghasilkan beberapa \omterm{def:b1:stereochemistry-in-depth:isomers}{isomer struktur} tetapi terutama
membentuk satu di antaranya adalah \emph{reaksi
regioselektif}\index{reaksi regioselektif}; reaksi yang dapat
menghasilkan beberapa \omterm{def:b1:stereochemistry-in-depth:isomers}{stereoisomer} tetapi terutama membentuk satu di
antaranya adalah
\emph{reaksi stereoselektif}\index{reaksi stereoselektif}.
\end{definition}

\begin{proposition}[Aturan Zaitsev]\label{prop:b1:elimination:zaitsev}
Jika beberapa hidrogen $\beta$ dapat dilepaskan, alkena utamanya biasanya
yang paling tersubstitusi (yang paling stabil), dan isomer \cip{E}
alih-alih isomer \cip{Z} (\emph{aturan Zaitsev}\index{aturan Zaitsev}).
Basa yang meruah, dan kendala geometris seperti syarat trans-diaksial,
dapat mengalahkan aturan itu.
\end{proposition}

\begin{proof}
Gugus-gugus alkil pada karbon-karbon ikatan rangkap dua menstabilkannya
(seperti gugus itu menstabilkan \omterm{def:b1:electronic-effects:intermediates}{karbokation}), dan alkena \cip{E}
menghindari berdesakannya gugus-gugus \textit{cis}. \omterm{def:b1:elementary-steps:energy-profile}{Keadaan transisi} E2
dan E1 mempunyai sebagian sifat ikatan rangkap dua, sehingga alkena yang
lebih stabil terbentuk lebih cepat. Basa yang meruah lebih mudah mencapai
hidrogen-hidrogen yang kurang terhalang; dan hidrogen yang tidak dapat
\omterm{def:b1:elimination:periplanar}{anti-periplanar} terhadap \omterm{def:b1:electronic-effects:nucleophile}{gugus pergi} sama sekali tidak dilepaskan oleh
E2, betapa pun stabilnya alkena yang akan dihasilkannya.
\end{proof}

\begin{method}[Meramalkan alkenanya]\label{met:b1:elimination:predict-alkene}
\begin{enumerate}
  \item Daftarkan karbon-karbon $\beta$ yang membawa hidrogen.
  \item Untuk E2, pertahankan hanya hidrogen yang dapat \omterm{def:b1:elimination:periplanar}{anti-periplanar}
        terhadap Y (pada cincin: trans-diaksial, dalam \omterm{def:b1:stereochemistry-in-depth:conformations}{kursi} yang
        benar-benar dapat terbentuk).
  \item Di antaranya, utamakan hidrogen yang menghasilkan alkena paling
        tersubstitusi (Zaitsev), kecuali basanya meruah.
  \item Untuk masing-masing, gambarlah \omterm{def:b1:stereochemistry-in-depth:representations}{proyeksi Newman} dalam \omterm{def:b1:stereochemistry-in-depth:conformations}{konformasi
        anti} untuk membaca geometri \cip{E}/\cip{Z}-nya.
\end{enumerate}
\end{method}

\section{Substitusi atau eliminasi?}

\begin{proposition}[Substitusi melawan eliminasi]\label{prop:b1:elimination:sn-vs-e}
Eliminasi didukung oleh \omterm{def:b1:predominant-reaction:strong-acid}{basa kuat} yang meruah, oleh suhu tinggi, dan oleh
\omterm{def:b1:nucleophilic-substitution:substitution}{substrat} tersier atau yang padat; substitusi oleh \omterm{def:b1:electronic-effects:nucleophile}{nukleofil} yang baik
yang merupakan \omterm{def:b1:predominant-reaction:strong-acid}{basa lemah}, oleh suhu rendah, dan oleh \omterm{def:b1:nucleophilic-substitution:substitution}{substrat} primer.
\end{proposition}

\begin{proof}
Kekuatan basa mendukung serangan pada hidrogen; keruahan menghalangi
serangan pada karbon yang padat tetapi tidak pada hidrogen yang terbuka.
\omterm{def:b1:electronic-effects:carbon-class}{Karbon tersier} tidak dapat dicapai SN2, dan \omterm{def:b1:electronic-effects:intermediates}{karbokationnya} mudah
kehilangan proton. Eliminasi mengubah satu molekul menjadi dua atau tiga
(alkena, \omterm{def:b1:electronic-effects:nucleophile}{gugus pergi}, basa yang terprotonasi): eliminasi lebih disukai
seiring naiknya suhu, sebuah hasil yang dibuktikan dalam jilid Tahun~2.
\end{proof}

\begin{omfigure}
\begin{tabular}{llll}
\hline
\omterm{def:b1:nucleophilic-substitution:substitution}{substrat} & \omterm{def:b1:predominant-reaction:strong-acid}{basa kuat}, kecil & \omterm{def:b1:predominant-reaction:strong-acid}{basa kuat}, meruah & \omterm{def:b1:predominant-reaction:strong-acid}{basa lemah}, \omterm{def:b1:electronic-effects:nucleophile}{nukleofil} baik \\
\hline
primer & SN2 (sebagian E2) & E2 & SN2 \\
sekunder & E2 (dan SN2) & E2 & SN2 (aprotik), SN1/E1 (protik) \\
tersier & E2 & E2 & SN1/E1 (protik) \\
\hline
\end{tabular}

\smallskip
{\small Orientasi awal: reaksi utama pada setiap kasus. Pemanasan
mendukung eliminasi pada setiap kolom.}
\end{omfigure}

\section{Latihan}

\begin{exercise}[$\star$]\label{exo:b1:elimination:1}
Tuliskan alkena-alkena yang dapat terbentuk melalui eliminasi HBr dari
2-bromobutana, dari 2-bromo-2-metilbutana, dan dari bromosikloheksana,
serta alkena yang diramalkan \omterm{prop:b1:elimination:zaitsev}{aturan Zaitsev} sebagai produk utama.
\end{exercise}

\begin{exercise}[$\star$]\label{exo:b1:elimination:2}
Tuliskan \omterm{def:b1:rate-laws:order}{hukum laju} E1 dan E2. Bagaimana keduanya dapat dibedakan secara
eksperimen untuk 2-bromo-2-metilpropana dalam etanol yang mengandung
etoksida?
\end{exercise}

\begin{exercise}[$\star$]\label{exo:b1:elimination:3}
Gambarlah \omterm{def:b1:stereochemistry-in-depth:representations}{proyeksi Newman} 2-bromobutana sepanjang C2--C3 dalam ketiga
\omterm{def:b1:stereochemistry-in-depth:conformations}{konformasi goyangnya}, lalu nyatakan mana yang dapat mengalami E2.
\end{exercise}

\begin{exercise}[$\star$]\label{exo:b1:elimination:4}
Ramalkan reaksi utamanya (SN1, SN2, E1, E2): bromoetana dengan natrium
hidroksida di dalam air pada \qty{25}{\celsius}; 2-bromo-2-metilpropana
dengan kalium tert-butoksida; 2-bromopropana dengan natrium iodida dalam
propanon; 2-bromo-2-metilpropana dalam etanol panas.
\end{exercise}

\begin{exercise}[$\star\star$]\label{exo:b1:elimination:5}
Tunjukkan, dengan \omterm{def:b1:stereochemistry-in-depth:representations}{proyeksi Newman}, bahwa eliminasi E2 HBr dari kedua
\omterm{def:b1:stereochemistry-in-depth:enantiomers}{diastereomer} 2-bromo-3-metilpentana yang dapat kehilangan hidrogen C3
menghasilkan kedua \omterm{def:b1:stereochemistry-in-depth:isomers}{stereoisomer} 3-metilpent-2-ena.
\end{exercise}

\begin{exercise}[$\star\star$]\label{exo:b1:elimination:6}
Tuliskan mekanisme dehidrasi 2-metilbutan-2-ol dalam asam sulfat pekat
yang panas (E1) dan tentukan alkena utamanya.
\end{exercise}

\begin{exercise}[$\star\star$]\label{exo:b1:elimination:7}
Jelaskan mengapa \textit{trans}-1-bromo-4-tert-butilsikloheksana bereaksi
sangat lambat melalui E2 sedangkan isomer \textit{cis}-nya bereaksi cepat.
(Gugus tert-butil tetap \omterm{def:b1:stereochemistry-in-depth:conformations}{ekuatorial}.)
\end{exercise}

\begin{exercise}[$\star\star$]\label{exo:b1:elimination:8}
2-Bromo-2-metilbutana dengan natrium etoksida terutama menghasilkan
2-metilbut-2-ena; dengan kalium tert-butoksida, terutama
2-metilbut-1-ena. Jelaskan.
\end{exercise}

\begin{exercise}[$\star\star$]\label{exo:b1:elimination:9}
Mengapa eliminasi tidak pernah teramati pada bromometana? Dan pada
1-bromo-2,2-dimetilpropana melalui E2?
\end{exercise}

\begin{exercise}[$\star\star\star$]\label{exo:b1:elimination:10}
Dalam etanol pada \qty{25}{\celsius}, 2-bromo-2-metilpropana menghasilkan
campuran eter (SN1) dan alkena (E1) yang perbandingannya tidak
bergantung pada konsentrasi \omterm{def:b1:nucleophilic-substitution:substitution}{substrat}. Tunjukkan bahwa kedua produk berasal
dari \omterm{def:b1:electronic-effects:intermediates}{karbokation} yang sama dan bahwa rasionya adalah rasio dua \omterm{def:b1:rate-laws:order}{tetapan
laju}.
\end{exercise}

\begin{exercise}[$\star\star\star$]\label{exo:b1:elimination:11}
Salah satu \omterm{def:b1:stereochemistry-in-depth:enantiomers}{diastereomer} 1,2-dibromo-1,2-difeniletana (yang meso)
menghasilkan, melalui E2 dengan etoksida, satu \omterm{def:b1:stereochemistry-in-depth:isomers}{stereoisomer}
1-bromo-1,2-difeniletena. Tentukan yang mana, dengan \omterm{def:b1:stereochemistry-in-depth:representations}{proyeksi Newman}.
\end{exercise}

\begin{exercise}[$\star\star\star$]\label{exo:b1:elimination:12}
Dengan energi-energi pada \cref{ch:b1:stereochemistry-in-depth}, jelaskan
mengapa reaksi E2 pada sikloheksana dapat diperlambat dengan faktor
ratusan jika \omterm{def:b1:electronic-effects:nucleophile}{gugus pergi} harus menjadi \omterm{def:b1:stereochemistry-in-depth:conformations}{aksial} dalam \omterm{def:b1:stereochemistry-in-depth:conformations}{kursi} yang padat.
Nyatakan lajunya sebagai hasil kali fraksi konformer dan sebuah \omterm{def:b1:rate-laws:order}{tetapan
laju}.
\end{exercise}

\section{Soal: Mentil Klorida dan Neomentil Klorida}

\begin{problem}[{Soal akhir pekan --- kursi-kursi dua klorida yang
merupakan diastereomer, hidrogen anti-periplanar yang tersedia pada
masing-masing, alkena yang terbentuk, dan rasio laju E2-nya}]
\label{pb:b1:elimination:1}
Mentil klorida dan neomentil klorida adalah klorida dari mentol dan
neomentol (\cref{ch:b1:stereochemistry-in-depth}): pada cincinnya, C1
membawa Cl, C2 gugus isopropil, C5 gugus metil. Pada mentil klorida
ketiga gugus itu dapat semuanya \omterm{def:b1:stereochemistry-in-depth:conformations}{ekuatorial}; pada neomentil klorida, Cl
\omterm{def:b1:stereochemistry-in-depth:conformations}{aksial} jika kedua gugus alkil \omterm{def:b1:stereochemistry-in-depth:conformations}{ekuatorial}. Keduanya direaksikan dengan
natrium etoksida dalam etanol, sebuah reaksi E2. Data soal:
\qty{95}{\percent} neomentil klorida berada dalam \omterm{def:b1:stereochemistry-in-depth:conformations}{kursi} dengan Cl
\omterm{def:b1:stereochemistry-in-depth:conformations}{aksial}; untuk mentil klorida, \omterm{def:b1:stereochemistry-in-depth:conformations}{kursi} dengan Cl \omterm{def:b1:stereochemistry-in-depth:conformations}{aksial} (ketiga gugus
\omterm{def:b1:stereochemistry-in-depth:conformations}{aksial}) merupakan fraksi \num{5.0e-3}; anggaplah setiap hidrogen
\omterm{def:b1:elimination:periplanar}{anti-periplanar} pada \omterm{def:b1:stereochemistry-in-depth:conformations}{kursi} yang reaktif bereaksi dengan \omterm{def:b1:rate-laws:order}{tetapan laju}
yang sama. Neomentil klorida menghasilkan \qty{78}{\percent} alkena
trisubstitusi.

\medskip
\noindent\textbf{Bagian I --- \omterm{def:b1:stereochemistry-in-depth:conformations}{Kursi-kursinya}.}
\begin{enumerate}
  \item Gambarlah \omterm{def:b1:stereochemistry-in-depth:conformations}{kursi} yang lebih disukai untuk mentil klorida.
  \item Gambarlah \omterm{def:b1:stereochemistry-in-depth:conformations}{kursi} yang lebih disukai untuk neomentil klorida.
  \item Apakah kedua klorida itu \omterm{def:b1:stereochemistry-in-depth:enantiomers}{enantiomer} atau \omterm{def:b1:stereochemistry-in-depth:enantiomers}{diastereomer}?
  \item Nyatakan syarat geometris E2 pada cincin.
  \item Pada \omterm{def:b1:stereochemistry-in-depth:conformations}{kursi} manakah mentil klorida dapat bereaksi? Gambarlah.
  \item Mengapa \omterm{def:b1:stereochemistry-in-depth:conformations}{kursi} itu jarang? Gunakan biaya energi gugus-gugus \omterm{def:b1:stereochemistry-in-depth:conformations}{aksial}.
\end{enumerate}

\medskip
\noindent\textbf{Bagian II --- Hidrogen-hidrogen anti.}
\begin{enumerate}[resume]
  \item Daftarkan karbon-karbon $\beta$ klorida itu (C2 dan C6) beserta
        hidrogennya.
  \item Pada \omterm{def:b1:stereochemistry-in-depth:conformations}{kursi} reaktif neomentil klorida, hidrogen manakah yang
        \omterm{def:b1:elimination:periplanar}{anti-periplanar} terhadap Cl?
  \item Pada \omterm{def:b1:stereochemistry-in-depth:conformations}{kursi} reaktif mentil klorida, apakah hidrogen C2 \omterm{def:b1:stereochemistry-in-depth:conformations}{aksial}?
        Apakah salah satu hidrogen C6 \omterm{def:b1:stereochemistry-in-depth:conformations}{aksial}?
  \item Berapa banyak hidrogen \omterm{def:b1:elimination:periplanar}{anti-periplanar} yang ditawarkan setiap
        klorida?
  \item Gambarlah \omterm{def:b1:stereochemistry-in-depth:representations}{proyeksi Newman} sepanjang C1--C6 untuk \omterm{def:b1:stereochemistry-in-depth:conformations}{kursi} reaktif
        mentil klorida.
  \item Mengapa eliminasi \omterm{def:b1:elimination:periplanar}{sin-periplanar} tidak dapat menyelamatkan reaksi
        itu?
\end{enumerate}

\medskip
\noindent\textbf{Bagian III --- Produk-produknya.}
\begin{enumerate}[resume]
  \item Namai alkena yang terbentuk dengan melepaskan hidrogen C2 dan
        alkena yang terbentuk dengan melepaskan salah satu hidrogen C6.
  \item Manakah yang didukung \omterm{prop:b1:elimination:zaitsev}{aturan Zaitsev}, dan mengapa?
  \item Apa yang dihasilkan neomentil klorida? Apakah rasio 78/22 sesuai
        dengan \omterm{prop:b1:elimination:zaitsev}{aturan Zaitsev}?
  \item Apa yang dihasilkan mentil klorida? Apakah hal ini sesuai dengan
        \omterm{prop:b1:elimination:zaitsev}{aturan Zaitsev}?
  \item Reaksi manakah yang stereospesifik, regioselektif, atau keduanya?
  \item Mengapa substitusi (SN2) hanya sedikit dengan etoksida di sini?
\end{enumerate}

\medskip
\noindent\textbf{Bagian IV --- Lajunya.}
\begin{enumerate}[resume]
  \item Tuliskan laju setiap reaksi sebagai fungsi fraksi \omterm{def:b1:stereochemistry-in-depth:conformations}{kursi} reaktif
        dan banyaknya hidrogen \omterm{def:b1:elimination:periplanar}{anti-periplanar}.
  \item Hitunglah laju relatif neomentil klorida.
  \item Hitunglah laju relatif mentil klorida.
  \item Hitunglah rasio laju $k(\text{neomentil})/k(\text{mentil})$.
\end{enumerate}
\end{problem}
